\section{Conclusión}\label{sec:conclusion}

El proyecto se propuso predecir la dificultad de pago del mes siguiente con modelos que, además de discriminar, explicaran sus decisiones. Se compararon un árbol, una regresión logística y un EBM sobre las mismas cinco particiones externas agrupadas, con validación anidada, selección por ROC-AUC y umbral por exactitud equilibrada. Esa homogeneidad permite comparar las familias sin que el reparto las favorezca.

El EBM obtuvo la mayor ROC-AUC ($0{,}7755 \pm 0{,}0077$) y la mejor exactitud equilibrada ($0{,}6997$ con umbral interno), seguido de la logística ($0{,}7515$) y del árbol ($0{,}7372$). La ventaja del EBM se repite en las cinco particiones, aunque el árbol conserva la mayor exactitud ($0{,}7793$) y precisión ($0{,}4805$). El umbral interno eleva la sensibilidad en más de $34$ puntos porcentuales, a costa de precisión y de exactitud, que cae por debajo de la referencia mayoritaria. No hay umbral neutro, porque fijarlo equivale a elegir qué error se tolera.

El modelo final fue un EBM sin interacciones, con \var{max_bins}=32, $10$ observaciones mínimas por hoja y umbral $0{,}20$. Aplicado al test, asigna $3\,208$ casos a la clase 0 y $1\,792$ a la clase 1, y se guarda en \texttt{outputs/test.csv} conservando filas y valores originales; sin etiquetas no pueden calcularse métricas.

La interpretabilidad no es una sola propiedad. Las reglas del árbol son inmediatas pero gruesas, los coeficientes de la logística compactos pero con referencias y escalas, y las funciones del EBM más fieles pero más costosas de inspeccionar. Una misma decisión no implica la misma explicación, y el umbral forma parte de ella.

Los límites acotan las conclusiones. Las cinco particiones proceden de entrenamientos solapados, no de réplicas independientes, así que su dispersión mide la sensibilidad al reparto. La familia se eligió con los resultados que se comparan, lo que puede introducir optimismo: el $0{,}7007$ final no acredita mejora externa sobre el $0{,}6997$. El test sin etiquetas impide afirmar generalización real, y la ventaja del EBM también depende de factores como la forma funcional o la penalización. Las mejoras posibles son la poda por coste-complejidad en el árbol, los splines o interacciones en la logística y la monotonía en variables como los retrasos en el EBM como se comentaba en la sección anterior.

El valor del proyecto reside en combinar rendimiento y explicaciones fieles sin confundir asociación con causa.
